\documentclass[10pt,a4paper]{amsart}
\usepackage[margin=19mm]{geometry}
\usepackage{amsmath,amssymb,mathtools}
\usepackage[scr=rsfs,cal=euler]{mathalfa}
\usepackage{xfrac}
\usepackage[unicode,hidelinks,bookmarks=false]{hyperref}
\newcommand{\ie}{i.e.\ }
\newcommand{\cf}{cf.\ }
\newcommand{\resp}{respectively\ }
\newcommand{\Subsection}{\subsection}
\providecommand{\nobreakdash}{\nobreak\hbox{-}}
\setlength{\emergencystretch}{2em}
\multlinegap0pt
\usepackage{bm}
\usepackage{bbm}
\usepackage{dsfont}
\usepackage{xspace}
\usepackage{cancel}
\usepackage{mathtools}
\usepackage{amsthm}
\makeatletter
\@ifundefined{theorem}{\newtheorem{theorem}{Theorem}}{}
\@ifundefined{lemma}{\newtheorem{lemma}[theorem]{Lemma}}{}
\@ifundefined{proposition}{\newtheorem{proposition}[theorem]{Proposition}}{}
\makeatother
\newcommand{\C}{\mathbb{C}}
\newcommand{\EC}{\mathcal{C}^\mathrm{EC}}
\title[On semi-finite vector bundles with connection over Kahler ma]{On semi-finite vector bundles with connection over Kahler manifolds}
\author{Sanjay Amrutiya, Indranil Biswas}
\date{Source: arXiv:2508.17048v1 (CC BY 4.0)}
\begin{document}
\maketitle

\noindent\emph{Source and licence.} On semi-finite vector bundles with connection over Kahler manifolds. Version arXiv:2508.17048v1 (CC BY 4.0), distributed under the Creative Commons Attribution 4.0 International licence, as recorded in the arXiv licence metadata for this exact version and shown on the abstract page https://arxiv.org/abs/2508.17048v1. Full rights notice: licenses/arxiv-2508.17048-CC-BY-4.0.txt.\par\smallskip

\noindent\textit{Editorial scope.} Counted unit: the Proposition whose source label is \texttt{prop-EC-rigid} in the article (source0.tex lines 343--345, source label \texttt{prop-EC-rigid}), delivered together with the article's own complete original proof of it (source0.tex lines 347--369, one \texttt{proof} environment of 23 lines). No mathematics is written, repaired or re-derived: statement and proof are the article's own text line for line. Required context retained uncounted: lemma-ext (lines 239--248); lem-t (lines 301--304). Every definition, display and label of those lines is retained unchanged. Omitted from this package and not redistributed: 1--238, 249--300, 305--342, 346--346, 370--525. Nothing inside a retained excerpt is omitted or altered: every retained body is delivered exactly as the article's lines are written, so no omission receipt of any kind accompanies this package. Citations: the retained excerpts carry no citation command at all, so no bibliography entry of the article is transcribed and no reference is redistributed. Reference adaptation: none was needed and none was made; every reference inside the delivered unit resolves inside it, so no unresolved reference or citation is produced. Printed slips kept literal: no printed word, symbol, hypothesis or number of the counted statement or of its proof was altered. Layout and class: the article's own class is \texttt{amsart} with packages including amsfonts, amssymb, hyperref, graphicx, latexsym, amsmath, amsthm, xy; the wrapper is the lane's pinned \texttt{amsart} editorial wrapper at 10pt on A4, which restates the theorem-like environments the retained text needs and adds the article's own macro definitions that the retained text needs (\texttt{\textbackslash C}, \texttt{\textbackslash EC}), with extra wrapper packages (bm, bbm, dsfont, xspace, cancel, mathtools). The delivered numbering is the unit's own. Assets: no retained span contains a figure environment, a graphics-inclusion command, a PDF-page inclusion, a table or a TikZ picture; no image file is copied into this package, the package declares no asset dependency and no image file is redistributed with it. Licence: version arXiv:2508.17048v1 of this article is distributed under the Creative Commons Attribution 4.0 International licence, as recorded in the arXiv licence metadata for this exact version and shown on the abstract page https://arxiv.org/abs/2508.17048v1. A full rights notice is delivered at licenses/arxiv-2508.17048-CC-BY-4.0.txt. Characters: every retained excerpt is pure ASCII, so the delivered engine, XeLaTeX with its default Latin Modern text font, reports no missing character.
\begin{lemma}\label{lemma-ext}
Let $0\,\longrightarrow\, (E_1,\, \nabla^{E_1})\,\stackrel{f}{\longrightarrow}\,
(E, \,\nabla^E)\, \stackrel{g}{\longrightarrow}\, (E_2,\, \nabla^{E_2})\,\longrightarrow\, 0$
be an exact sequence in $\mathbf{Conn}(X)$, meaning the underlying sequence of holomorphic 
vector bundles
is exact with $\nabla^E\big\vert_{E_1} \,=\, \nabla^{E_1}$ and $\nabla^{E_2}$ coincides
with the connection on $E_2$ induced by $\nabla^E$.
If $(E_1,\, \nabla^{E_1})$ and $(E_2, \nabla^{E_2})$ are objects of $\mathrm{\EC(X)}$, 
then $(E, \,\nabla^E)$ is also an object of $\EC(X)$.
\end{lemma}

\begin{lemma}\label{lem-t}
For any object $(E,\, \nabla^E)$ of $\EC(X)$, the dual vector bundle $E^*$ equipped with the
connection induced by $\nabla^E$ is also an object of $\EC(X)$.
\end{lemma}

\begin{proposition}\label{prop-EC-rigid}
The full subcategory $\EC(X)$ of $\mathbf{Conn}(X)$ is a $\C$-linear rigid tensor subcategory. 
\end{proposition}

\begin{proof}
We will show that $\EC(X)$ is a tensor category. Take two objects $(E,\, \nabla^E)$ and
$(F,\, \nabla^F)$ in $\EC(X)$. We need to show that $E \otimes F$ together with the induced 
connection $\nabla^{E\otimes F})$ is an object of $\EC(X)$. 
Since $(E,\, \nabla^E)$ is an object in $\EC(X)$, there exist a filtration
\begin{equation}\label{eq:2}
E\,=\,E_0\,\supseteq\, E_1\, \supseteq\, E_2\, \supseteq\, \cdots\, \supseteq\, E_n\,\supseteq\, 
E_{n+1}\, =\, 0
\end{equation}
with $\nabla^E(E_i)\, \subseteq\, E_i \otimes_{\mathcal{O}_X} \Omega_X^1$, and the flat holomorphic
connection on each quotient $E_i/E_{i+1}$, $0\, \leq\, i\, \leq\, n$, induced by $\nabla^E$ 
has finite monodromy group. Tensoring \eqref{eq:2} with $F$, we get a filtration 
$$
E\otimes F\,=\,E_0\otimes F\,\supseteq\, E_1\otimes F\, \supseteq\, \cdots\, \supseteq\,
E_n\otimes F\,\supseteq\, E_{n+1}\otimes F\, =\, 0
$$ 
with $(E_i \otimes F)/(E_{i+1} \otimes F)\, \simeq\, (E_i/E_{i+1})\otimes F$ and it is compatible 
with the induced connection $\nabla^{E\otimes F}$. By Lemma \ref{lemma-ext}, the 
category $\EC(X)$ is closed under taking extension in $\mathbf{Conn}(X)$. Consequently, it follows 
that $(E \otimes F,\, \nabla^{E\otimes F})$ is an object of $\EC(X)$. 

In Lemma \ref{lem-t}, it was shown that $\EC(X)$ is closed under taking duals. This completes the proof.
\end{proof}

\end{document}
